%%%%%%%%%%%%%%%%%%%%%%%%%%%%%%%%%%%%%%%%%%%%%%%%%%%%%%%%%%%%
% 8.8 ANALYSIS OF VARIANCE
% The Composition of Advanced Mathematics
%%%%%%%%%%%%%%%%%%%%%%%%%%%%%%%%%%%%%%%%%%%%%%%%%%%%%%%%%%%%

\section{Analysis of Variance}
\label{sec:analysis-of-variance}

\prerequisite{
Descriptive statistics, sampling distributions, statistical estimation,
confidence intervals, hypothesis testing, regression, normal
distributions, chi-square and \(F\)-distributions, sums of squares,
linear algebra, and basic experimental design concepts.
}

\learningobjectives{
By the end of this section, the reader should be able to:
\begin{itemize}
    \item explain the purpose of analysis of variance;
    \item distinguish within-group and between-group variability;
    \item formulate a one-way ANOVA model;
    \item state the null and alternative hypotheses for one-way ANOVA;
    \item compute group means and the grand mean;
    \item derive the ANOVA sums-of-squares decomposition;
    \item compute treatment and error sums of squares;
    \item compute mean squares;
    \item construct the one-way ANOVA \(F\)-statistic;
    \item interpret ANOVA tables;
    \item state the classical ANOVA assumptions;
    \item understand ANOVA as a regression model;
    \item understand indicator-variable representations of group effects;
    \item understand fixed-effects and random-effects ideas conceptually;
    \item distinguish omnibus and pairwise hypotheses;
    \item perform planned contrasts;
    \item understand orthogonal contrasts;
    \item understand multiple-comparison problems;
    \item apply Bonferroni adjustments conceptually;
    \item understand Tukey simultaneous comparisons conceptually;
    \item understand Welch's ANOVA for unequal variances;
    \item understand nonparametric alternatives conceptually;
    \item understand two-way ANOVA;
    \item distinguish main effects and interaction effects;
    \item understand balanced factorial designs;
    \item understand repeated-measures ANOVA conceptually;
    \item understand blocking through ANOVA;
    \item understand nested factors conceptually;
    \item understand ANCOVA conceptually;
    \item interpret effect sizes such as \(\eta^2\) and partial
          \(\eta^2\);
    \item diagnose ANOVA assumptions using residuals;
    \item connect ANOVA with regression and experimental design.
\end{itemize}
}

%%%%%%%%%%%%%%%%%%%%%%%%%%%%%%%%%%%%%%%%%%%%%%%%%%%%%%%%%%%%
\subsection{What Is Analysis of Variance?}
\label{subsec:what-is-anova}
%%%%%%%%%%%%%%%%%%%%%%%%%%%%%%%%%%%%%%%%%%%%%%%%%%%%%%%%%%%%

Analysis of variance, abbreviated ANOVA, is a framework for comparing
population means across multiple groups.

Suppose there are \(k\) populations with means

\[
\mu_1,
\mu_2,
\ldots,
\mu_k.
\]

The basic one-way ANOVA null hypothesis is

\[
\boxed{
H_0:
\mu_1=\mu_2=\cdots=\mu_k.
}
\]

The alternative is

\[
\boxed{
H_A:
\text{not all population means are equal}.
}
\]

ANOVA compares two sources of variability:

\[
\boxed{
\text{variation between groups}
}
\]

and

\[
\boxed{
\text{variation within groups}.
}
\]

%%%%%%%%%%%%%%%%%%%%%%%%%%%%%%%%%%%%%%%%%%%%%%%%%%%%%%%%%%%%
\subsection{Why Compare Variances to Test Means?}
\label{subsec:why-anova-compares-variance}
%%%%%%%%%%%%%%%%%%%%%%%%%%%%%%%%%%%%%%%%%%%%%%%%%%%%%%%%%%%%

The name ``analysis of variance'' may seem surprising because the primary
hypothesis concerns means.

The idea is that if all group means are truly equal, then differences
among observed sample means should be explainable by ordinary
within-group variation.

If observed group means are much farther apart than expected relative to
within-group variability, this provides evidence against

\[
H_0.
\]

Thus ANOVA tests means by comparing variance-like quantities.

%%%%%%%%%%%%%%%%%%%%%%%%%%%%%%%%%%%%%%%%%%%%%%%%%%%%%%%%%%%%
\subsection{One-Way ANOVA Setting}
\label{subsec:one-way-anova-setting}
%%%%%%%%%%%%%%%%%%%%%%%%%%%%%%%%%%%%%%%%%%%%%%%%%%%%%%%%%%%%

Suppose one categorical factor has

\[
k
\]

levels.

Let

\[
Y_{ij}
\]

denote observation \(j\) in group \(i\), where

\[
i=1,\ldots,k
\]

and

\[
j=1,\ldots,n_i.
\]

The total sample size is

\[
\boxed{
N
=
\sum_{i=1}^{k}n_i.
}
\]

%%%%%%%%%%%%%%%%%%%%%%%%%%%%%%%%%%%%%%%%%%%%%%%%%%%%%%%%%%%%
\subsection{One-Way ANOVA Model}
\label{subsec:one-way-anova-model}
%%%%%%%%%%%%%%%%%%%%%%%%%%%%%%%%%%%%%%%%%%%%%%%%%%%%%%%%%%%%

A common model is

\[
\boxed{
Y_{ij}
=
\mu_i
+
\varepsilon_{ij},
}
\]

where

\[
\mu_i
\]

is the population mean of group \(i\).

An equivalent effects model is

\[
\boxed{
Y_{ij}
=
\mu
+
\tau_i
+
\varepsilon_{ij},
}
\]

where

\[
\mu
\]

is an overall mean and

\[
\tau_i
\]

is the effect of group \(i\).

%%%%%%%%%%%%%%%%%%%%%%%%%%%%%%%%%%%%%%%%%%%%%%%%%%%%%%%%%%%%
\subsection{Identifiability Constraint}
\label{subsec:anova-identifiability}
%%%%%%%%%%%%%%%%%%%%%%%%%%%%%%%%%%%%%%%%%%%%%%%%%%%%%%%%%%%%

The decomposition

\[
\mu+\tau_i
\]

is not unique unless a constraint is imposed.

A common weighted constraint is

\[
\boxed{
\sum_{i=1}^{k}
n_i\tau_i
=
0.
}
\]

For a balanced design with

\[
n_i=n,
\]

this reduces to

\[
\boxed{
\sum_{i=1}^{k}
\tau_i
=
0.
}
\]

Another common approach is to designate one group as a reference group.

%%%%%%%%%%%%%%%%%%%%%%%%%%%%%%%%%%%%%%%%%%%%%%%%%%%%%%%%%%%%
\subsection{Group Means}
\label{subsec:anova-group-means}
%%%%%%%%%%%%%%%%%%%%%%%%%%%%%%%%%%%%%%%%%%%%%%%%%%%%%%%%%%%%

For group \(i\), the sample mean is

\[
\boxed{
\overline{Y}_{i\cdot}
=
\frac1{n_i}
\sum_{j=1}^{n_i}
Y_{ij}.
}
\]

The dot indicates averaging over the \(j\)-index.

For observed data,

\[
\overline{y}_{i\cdot}
\]

is the corresponding numerical group mean.

%%%%%%%%%%%%%%%%%%%%%%%%%%%%%%%%%%%%%%%%%%%%%%%%%%%%%%%%%%%%
\subsection{Grand Mean}
\label{subsec:anova-grand-mean}
%%%%%%%%%%%%%%%%%%%%%%%%%%%%%%%%%%%%%%%%%%%%%%%%%%%%%%%%%%%%

The grand mean is the mean of all \(N\) observations:

\[
\boxed{
\overline{Y}_{\cdot\cdot}
=
\frac1N
\sum_{i=1}^{k}
\sum_{j=1}^{n_i}
Y_{ij}.
}
\]

Equivalently,

\[
\boxed{
\overline{Y}_{\cdot\cdot}
=
\frac{
\sum_{i=1}^{k}
n_i\overline{Y}_{i\cdot}
}{
N
}.
}
\]

Thus the grand mean is a weighted average of the group means.

%%%%%%%%%%%%%%%%%%%%%%%%%%%%%%%%%%%%%%%%%%%%%%%%%%%%%%%%%%%%
\subsection{Worked Example: Grand Mean}
\label{subsec:worked-anova-grand-mean}
%%%%%%%%%%%%%%%%%%%%%%%%%%%%%%%%%%%%%%%%%%%%%%%%%%%%%%%%%%%%

\begin{workedexamplebox}{Grand Mean from Unequal Group Sizes}

Suppose three groups have

\[
n_1=4,
\qquad
n_2=3,
\qquad
n_3=5,
\]

with means

\[
\overline{y}_{1\cdot}=10,
\qquad
\overline{y}_{2\cdot}=12,
\qquad
\overline{y}_{3\cdot}=8.
\]

Then

\[
N=12.
\]

The grand mean is

\[
\overline{y}_{\cdot\cdot}
=
\frac{
4(10)+3(12)+5(8)
}{
12
}.
\]

Therefore,

\begin{answerbox}
\[
\boxed{
\overline{y}_{\cdot\cdot}
=
\frac{116}{12}
\approx
9.667.
}
\]
\end{answerbox}

\end{workedexamplebox}

%%%%%%%%%%%%%%%%%%%%%%%%%%%%%%%%%%%%%%%%%%%%%%%%%%%%%%%%%%%%
\subsection{Total Variation}
\label{subsec:anova-total-variation}
%%%%%%%%%%%%%%%%%%%%%%%%%%%%%%%%%%%%%%%%%%%%%%%%%%%%%%%%%%%%

The total sum of squares is

\[
\boxed{
SS_T
=
\sum_{i=1}^{k}
\sum_{j=1}^{n_i}
\left(
Y_{ij}
-
\overline{Y}_{\cdot\cdot}
\right)^2.
}
\]

This measures total observed variability around the grand mean.

Its degrees of freedom are

\[
\boxed{
N-1.
}
\]

%%%%%%%%%%%%%%%%%%%%%%%%%%%%%%%%%%%%%%%%%%%%%%%%%%%%%%%%%%%%
\subsection{Between-Group Variation}
\label{subsec:between-group-variation}
%%%%%%%%%%%%%%%%%%%%%%%%%%%%%%%%%%%%%%%%%%%%%%%%%%%%%%%%%%%%

The between-group or treatment sum of squares is

\[
\boxed{
SS_A
=
\sum_{i=1}^{k}
n_i
\left(
\overline{Y}_{i\cdot}
-
\overline{Y}_{\cdot\cdot}
\right)^2.
}
\]

Alternative notations include

\[
SS_{\text{Between}},
\qquad
SS_{\text{Treatment}},
\]

or

\[
SS_{\text{Model}}.
\]

This measures variability among group means.

%%%%%%%%%%%%%%%%%%%%%%%%%%%%%%%%%%%%%%%%%%%%%%%%%%%%%%%%%%%%
\subsection{Within-Group Variation}
\label{subsec:within-group-variation}
%%%%%%%%%%%%%%%%%%%%%%%%%%%%%%%%%%%%%%%%%%%%%%%%%%%%%%%%%%%%

The within-group or error sum of squares is

\[
\boxed{
SS_E
=
\sum_{i=1}^{k}
\sum_{j=1}^{n_i}
\left(
Y_{ij}
-
\overline{Y}_{i\cdot}
\right)^2.
}
\]

Alternative notation includes

\[
SS_{\text{Within}}
\]

or

\[
SS_{\text{Error}}.
\]

This measures variation of observations around their own group means.

%%%%%%%%%%%%%%%%%%%%%%%%%%%%%%%%%%%%%%%%%%%%%%%%%%%%%%%%%%%%
\subsection{ANOVA Sum-of-Squares Identity}
\label{subsec:anova-sum-squares-identity}
%%%%%%%%%%%%%%%%%%%%%%%%%%%%%%%%%%%%%%%%%%%%%%%%%%%%%%%%%%%%

A central identity is

\[
\boxed{
SS_T
=
SS_A
+
SS_E.
}
\]

Thus

\[
\boxed{
\text{total variation}
=
\text{between-group variation}
+
\text{within-group variation}.
}
\]

This is the ANOVA decomposition.

%%%%%%%%%%%%%%%%%%%%%%%%%%%%%%%%%%%%%%%%%%%%%%%%%%%%%%%%%%%%
\subsection{Deriving the ANOVA Decomposition}
\label{subsec:derive-anova-decomposition}
%%%%%%%%%%%%%%%%%%%%%%%%%%%%%%%%%%%%%%%%%%%%%%%%%%%%%%%%%%%%

For every observation,

\[
Y_{ij}
-
\overline{Y}_{\cdot\cdot}
=
\left(
Y_{ij}
-
\overline{Y}_{i\cdot}
\right)
+
\left(
\overline{Y}_{i\cdot}
-
\overline{Y}_{\cdot\cdot}
\right).
\]

Squaring,

\[
\begin{aligned}
\left(
Y_{ij}
-
\overline{Y}_{\cdot\cdot}
\right)^2
&=
\left(
Y_{ij}
-
\overline{Y}_{i\cdot}
\right)^2
\\
&\quad+
2
\left(
Y_{ij}
-
\overline{Y}_{i\cdot}
\right)
\left(
\overline{Y}_{i\cdot}
-
\overline{Y}_{\cdot\cdot}
\right)
\\
&\quad+
\left(
\overline{Y}_{i\cdot}
-
\overline{Y}_{\cdot\cdot}
\right)^2.
\end{aligned}
\]

Summing over \(j\), the cross term vanishes because

\[
\sum_{j=1}^{n_i}
\left(
Y_{ij}
-
\overline{Y}_{i\cdot}
\right)
=
0.
\]

Therefore,

\[
\boxed{
SS_T
=
SS_A+SS_E.
}
\]

%%%%%%%%%%%%%%%%%%%%%%%%%%%%%%%%%%%%%%%%%%%%%%%%%%%%%%%%%%%%
\subsection{Degrees of Freedom}
\label{subsec:anova-degrees-freedom}
%%%%%%%%%%%%%%%%%%%%%%%%%%%%%%%%%%%%%%%%%%%%%%%%%%%%%%%%%%%%

The total degrees of freedom are

\[
\boxed{
N-1.
}
\]

The between-group degrees of freedom are

\[
\boxed{
k-1.
}
\]

The within-group degrees of freedom are

\[
\boxed{
N-k.
}
\]

These satisfy

\[
\boxed{
N-1
=
(k-1)
+
(N-k).
}
\]

%%%%%%%%%%%%%%%%%%%%%%%%%%%%%%%%%%%%%%%%%%%%%%%%%%%%%%%%%%%%
\subsection{Mean Square Between Groups}
\label{subsec:mean-square-between}
%%%%%%%%%%%%%%%%%%%%%%%%%%%%%%%%%%%%%%%%%%%%%%%%%%%%%%%%%%%%

The treatment mean square is

\[
\boxed{
MS_A
=
\frac{
SS_A
}{
k-1
}.
}
\]

This is a variance-like measure of variation among group means.

Under the null hypothesis, it estimates the common error variance.

Under alternatives, it also reflects differences among population means.

%%%%%%%%%%%%%%%%%%%%%%%%%%%%%%%%%%%%%%%%%%%%%%%%%%%%%%%%%%%%
\subsection{Mean Square Error}
\label{subsec:anova-mean-square-error}
%%%%%%%%%%%%%%%%%%%%%%%%%%%%%%%%%%%%%%%%%%%%%%%%%%%%%%%%%%%%

The error mean square is

\[
\boxed{
MS_E
=
\frac{
SS_E
}{
N-k
}.
}
\]

Under the equal-variance ANOVA model,

\[
MS_E
\]

estimates the common error variance

\[
\boxed{
\sigma^2.
}
\]

Thus

\[
\boxed{
\widehat{\sigma}^2
=
MS_E.
}
\]

%%%%%%%%%%%%%%%%%%%%%%%%%%%%%%%%%%%%%%%%%%%%%%%%%%%%%%%%%%%%
\subsection{The ANOVA \(F\)-Statistic}
\label{subsec:anova-f-statistic}
%%%%%%%%%%%%%%%%%%%%%%%%%%%%%%%%%%%%%%%%%%%%%%%%%%%%%%%%%%%%

The one-way ANOVA test statistic is

\[
\boxed{
F
=
\frac{
MS_A
}{
MS_E
}.
}
\]

Under

\[
H_0:
\mu_1=\cdots=\mu_k
\]

and the classical normal model,

\[
\boxed{
F
\sim
F_{k-1,N-k}.
}
\]

Large values of \(F\) provide evidence against the null.

%%%%%%%%%%%%%%%%%%%%%%%%%%%%%%%%%%%%%%%%%%%%%%%%%%%%%%%%%%%%
\subsection{Why Large \(F\) Values Matter}
\label{subsec:why-large-anova-f}
%%%%%%%%%%%%%%%%%%%%%%%%%%%%%%%%%%%%%%%%%%%%%%%%%%%%%%%%%%%%

Under \(H_0\), both

\[
MS_A
\]

and

\[
MS_E
\]

estimate approximately the same error variance.

Therefore one expects

\[
F
\]

to be reasonably near \(1\).

If population means differ substantially,

\[
MS_A
\]

tends to become larger.

Thus

\[
\boxed{
F\gg1
}
\]

suggests meaningful between-group variation relative to within-group
noise.

%%%%%%%%%%%%%%%%%%%%%%%%%%%%%%%%%%%%%%%%%%%%%%%%%%%%%%%%%%%%
\subsection{ANOVA Table}
\label{subsec:one-way-anova-table}
%%%%%%%%%%%%%%%%%%%%%%%%%%%%%%%%%%%%%%%%%%%%%%%%%%%%%%%%%%%%

The standard one-way ANOVA table is

\[
\begin{array}{c|c|c|c|c}
\text{Source}
&
\text{df}
&
\text{SS}
&
\text{MS}
&
F
\\
\hline
\text{Groups}
&
k-1
&
SS_A
&
MS_A
&
MS_A/MS_E
\\
\text{Error}
&
N-k
&
SS_E
&
MS_E
&
\\
\text{Total}
&
N-1
&
SS_T
&
&
\end{array}
\]

This format summarizes the entire one-way ANOVA calculation.

%%%%%%%%%%%%%%%%%%%%%%%%%%%%%%%%%%%%%%%%%%%%%%%%%%%%%%%%%%%%
\subsection{Worked Example: One-Way ANOVA Table}
\label{subsec:worked-one-way-anova-table}
%%%%%%%%%%%%%%%%%%%%%%%%%%%%%%%%%%%%%%%%%%%%%%%%%%%%%%%%%%%%

\begin{workedexamplebox}{Constructing an ANOVA \(F\)-Statistic}

Suppose

\[
k=3,
\qquad
N=18,
\]

and the sums of squares are

\[
SS_A=120,
\qquad
SS_E=150.
\]

The between-group degrees of freedom are

\[
3-1=2.
\]

The error degrees of freedom are

\[
18-3=15.
\]

Thus,

\[
MS_A
=
\frac{120}{2}
=
60,
\]

and

\[
MS_E
=
\frac{150}{15}
=
10.
\]

Therefore,

\begin{answerbox}
\[
\boxed{
F
=
\frac{60}{10}
=
6.
}
\]

with degrees of freedom

\[
\boxed{
(2,15).
}
\]
\end{answerbox}

\end{workedexamplebox}

%%%%%%%%%%%%%%%%%%%%%%%%%%%%%%%%%%%%%%%%%%%%%%%%%%%%%%%%%%%%
\subsection{One-Way ANOVA Hypothesis Decision}
\label{subsec:one-way-anova-decision}
%%%%%%%%%%%%%%%%%%%%%%%%%%%%%%%%%%%%%%%%%%%%%%%%%%%%%%%%%%%%

For a significance level \(\alpha\), reject \(H_0\) when the observed
\(F\)-statistic lies sufficiently far in the upper tail of

\[
F_{k-1,N-k}.
\]

Equivalently, reject when

\[
\boxed{
p\text{-value}
\leq
\alpha.
}
\]

The ANOVA \(F\)-test is an upper-tailed test because large ratios indicate
excess between-group variation.

%%%%%%%%%%%%%%%%%%%%%%%%%%%%%%%%%%%%%%%%%%%%%%%%%%%%%%%%%%%%
\subsection{What ANOVA Rejection Means}
\label{subsec:what-anova-rejection-means}
%%%%%%%%%%%%%%%%%%%%%%%%%%%%%%%%%%%%%%%%%%%%%%%%%%%%%%%%%%%%

If the ANOVA null is rejected, the conclusion is

\[
\boxed{
\text{not all means are equal}.
}
\]

It does not immediately tell us:

\[
\boxed{
\text{which means differ},
}
\]

or

\[
\boxed{
\text{how large the differences are}.
}
\]

Additional contrasts or multiple-comparison procedures are required.

%%%%%%%%%%%%%%%%%%%%%%%%%%%%%%%%%%%%%%%%%%%%%%%%%%%%%%%%%%%%
\subsection{What ANOVA Failure to Reject Means}
\label{subsec:anova-fail-reject}
%%%%%%%%%%%%%%%%%%%%%%%%%%%%%%%%%%%%%%%%%%%%%%%%%%%%%%%%%%%%

If the ANOVA test is not significant, one should not conclude that all
means are exactly equal.

A nonsignificant result may reflect:

\[
\boxed{
\text{small true differences},
}
\]

\[
\boxed{
\text{high within-group variability},
}
\]

or

\[
\boxed{
\text{insufficient sample size}.
}
\]

Thus:

\[
\boxed{
\text{fail to reject equality}
\neq
\text{prove equality}.
}
\]

%%%%%%%%%%%%%%%%%%%%%%%%%%%%%%%%%%%%%%%%%%%%%%%%%%%%%%%%%%%%
\subsection{Classical One-Way ANOVA Assumptions}
\label{subsec:anova-assumptions}
%%%%%%%%%%%%%%%%%%%%%%%%%%%%%%%%%%%%%%%%%%%%%%%%%%%%%%%%%%%%

The classical one-way fixed-effects ANOVA commonly assumes:

\[
\boxed{
\text{independent observations},
}
\]

\[
\boxed{
\text{normal errors within groups},
}
\]

and

\[
\boxed{
\text{equal population variances}.
}
\]

That is,

\[
\boxed{
\varepsilon_{ij}
\overset{\text{ind.}}{\sim}
N(0,\sigma^2).
}
\]

%%%%%%%%%%%%%%%%%%%%%%%%%%%%%%%%%%%%%%%%%%%%%%%%%%%%%%%%%%%%
\subsection{Independence Assumption}
\label{subsec:anova-independence}
%%%%%%%%%%%%%%%%%%%%%%%%%%%%%%%%%%%%%%%%%%%%%%%%%%%%%%%%%%%%

Independence is primarily a design assumption.

It may be justified through:

\[
\boxed{
\text{independent random sampling},
}
\]

or

\[
\boxed{
\text{independent random assignment}.
}
\]

Repeated measurements on the same individual or observations from the
same cluster generally violate simple independence.

%%%%%%%%%%%%%%%%%%%%%%%%%%%%%%%%%%%%%%%%%%%%%%%%%%%%%%%%%%%%
\subsection{Normality Assumption}
\label{subsec:anova-normality}
%%%%%%%%%%%%%%%%%%%%%%%%%%%%%%%%%%%%%%%%%%%%%%%%%%%%%%%%%%%%

Normality concerns the distribution of errors within groups.

ANOVA is often reasonably robust to moderate nonnormality when sample
sizes are not extremely small and groups are reasonably balanced.

However, strong skewness, heavy tails, or severe outliers can affect the
test substantially.

Residual diagnostics should therefore be examined.

%%%%%%%%%%%%%%%%%%%%%%%%%%%%%%%%%%%%%%%%%%%%%%%%%%%%%%%%%%%%
\subsection{Equal-Variance Assumption}
\label{subsec:anova-equal-variance}
%%%%%%%%%%%%%%%%%%%%%%%%%%%%%%%%%%%%%%%%%%%%%%%%%%%%%%%%%%%%

Classical ANOVA assumes

\[
\boxed{
\sigma_1^2
=
\sigma_2^2
=
\cdots
=
\sigma_k^2
=
\sigma^2.
}
\]

This is called homogeneity of variance or homoskedasticity.

Severe variance differences can distort the ordinary ANOVA test,
especially when group sizes are unequal.

%%%%%%%%%%%%%%%%%%%%%%%%%%%%%%%%%%%%%%%%%%%%%%%%%%%%%%%%%%%%
\subsection{Residuals in One-Way ANOVA}
\label{subsec:anova-residuals}
%%%%%%%%%%%%%%%%%%%%%%%%%%%%%%%%%%%%%%%%%%%%%%%%%%%%%%%%%%%%

For observation \(Y_{ij}\), the fitted value is

\[
\boxed{
\widehat{Y}_{ij}
=
\overline{Y}_{i\cdot}.
}
\]

Therefore the residual is

\[
\boxed{
e_{ij}
=
Y_{ij}
-
\overline{Y}_{i\cdot}.
}
\]

Thus one-way ANOVA residuals are simply deviations from group means.

%%%%%%%%%%%%%%%%%%%%%%%%%%%%%%%%%%%%%%%%%%%%%%%%%%%%%%%%%%%%
\subsection{ANOVA Residual Diagnostics}
\label{subsec:anova-residual-diagnostics}
%%%%%%%%%%%%%%%%%%%%%%%%%%%%%%%%%%%%%%%%%%%%%%%%%%%%%%%%%%%%

Useful diagnostics include:

\[
\boxed{
\text{residuals versus fitted values},
}
\]

\[
\boxed{
\text{residuals versus group},
}
\]

and

\[
\boxed{
\text{normal Q--Q plots}.
}
\]

A funnel pattern may indicate unequal variance.

Extreme residuals may indicate outliers or model misspecification.

%%%%%%%%%%%%%%%%%%%%%%%%%%%%%%%%%%%%%%%%%%%%%%%%%%%%%%%%%%%%
\subsection{ANOVA as Regression}
\label{subsec:anova-as-regression}
%%%%%%%%%%%%%%%%%%%%%%%%%%%%%%%%%%%%%%%%%%%%%%%%%%%%%%%%%%%%

One-way ANOVA is a special case of linear regression.

Suppose there are \(k\) groups.

Choose group \(1\) as the reference group and define indicator variables

\[
D_2,\ldots,D_k.
\]

Then

\[
\boxed{
Y
=
\beta_0
+
\beta_2D_2
+
\cdots
+
\beta_kD_k
+
\varepsilon.
}
\]

Here,

\[
\beta_0
=
\mu_1,
\]

and

\[
\boxed{
\beta_j
=
\mu_j-\mu_1.
}
\]

%%%%%%%%%%%%%%%%%%%%%%%%%%%%%%%%%%%%%%%%%%%%%%%%%%%%%%%%%%%%
\subsection{ANOVA Null Hypothesis in Regression Form}
\label{subsec:anova-null-regression-form}
%%%%%%%%%%%%%%%%%%%%%%%%%%%%%%%%%%%%%%%%%%%%%%%%%%%%%%%%%%%%

The ANOVA null

\[
\mu_1=\cdots=\mu_k
\]

is equivalent to

\[
\boxed{
H_0:
\beta_2
=
\cdots
=
\beta_k
=
0.
}
\]

Thus the one-way ANOVA \(F\)-test is an overall regression \(F\)-test for
the group indicators.

This establishes the formal connection

\[
\boxed{
\text{ANOVA}
\subset
\text{linear models}.
}
\]

%%%%%%%%%%%%%%%%%%%%%%%%%%%%%%%%%%%%%%%%%%%%%%%%%%%%%%%%%%%%
\subsection{Two-Group ANOVA and the \(t\)-Test}
\label{subsec:two-group-anova-t-test}
%%%%%%%%%%%%%%%%%%%%%%%%%%%%%%%%%%%%%%%%%%%%%%%%%%%%%%%%%%%%

When

\[
k=2,
\]

the ordinary equal-variance one-way ANOVA \(F\)-test is equivalent to the
pooled two-sample \(t\)-test.

Specifically,

\[
\boxed{
F=t^2.
}
\]

Thus ANOVA generalizes the equal-variance two-sample comparison to more
than two groups.

%%%%%%%%%%%%%%%%%%%%%%%%%%%%%%%%%%%%%%%%%%%%%%%%%%%%%%%%%%%%
\subsection{Worked Example: \(F=t^2\)}
\label{subsec:worked-f-equals-t-square}
%%%%%%%%%%%%%%%%%%%%%%%%%%%%%%%%%%%%%%%%%%%%%%%%%%%%%%%%%%%%

\begin{workedexamplebox}{Two Groups}

Suppose a pooled two-sample comparison produces

\[
t=2.5.
\]

The equivalent one-way ANOVA statistic is

\[
F=t^2.
\]

Therefore,

\begin{answerbox}
\[
\boxed{
F
=
(2.5)^2
=
6.25.
}
\]
\end{answerbox}

The two procedures yield the same two-sided significance conclusion under
the common assumptions.

\end{workedexamplebox}

%%%%%%%%%%%%%%%%%%%%%%%%%%%%%%%%%%%%%%%%%%%%%%%%%%%%%%%%%%%%
\subsection{Planned Contrasts}
\label{subsec:planned-contrasts}
%%%%%%%%%%%%%%%%%%%%%%%%%%%%%%%%%%%%%%%%%%%%%%%%%%%%%%%%%%%%

ANOVA asks an omnibus question.

A contrast asks a more specific question about group means.

A linear contrast has the form

\[
\boxed{
L
=
\sum_{i=1}^{k}
c_i\mu_i,
}
\]

where the coefficients satisfy

\[
\boxed{
\sum_{i=1}^{k}
c_i=0.
}
\]

The sample estimator is

\[
\boxed{
\widehat{L}
=
\sum_{i=1}^{k}
c_i\overline{Y}_{i\cdot}.
}
\]

%%%%%%%%%%%%%%%%%%%%%%%%%%%%%%%%%%%%%%%%%%%%%%%%%%%%%%%%%%%%
\subsection{Worked Example: Contrast}
\label{subsec:worked-anova-contrast}
%%%%%%%%%%%%%%%%%%%%%%%%%%%%%%%%%%%%%%%%%%%%%%%%%%%%%%%%%%%%

\begin{workedexamplebox}{Comparing One Group to the Average of Two Others}

Suppose three means are

\[
\mu_1,
\mu_2,
\mu_3.
\]

To compare group \(1\) with the average of groups \(2\) and \(3\), use

\[
L
=
\mu_1
-
\frac12\mu_2
-
\frac12\mu_3.
\]

The coefficients are

\[
1,
\quad
-\frac12,
\quad
-\frac12.
\]

Their sum is

\[
1-\frac12-\frac12=0.
\]

\begin{answerbox}
\[
\boxed{
L
=
\mu_1
-
\frac{
\mu_2+\mu_3
}{2}.
}
\]
\end{answerbox}

\end{workedexamplebox}

%%%%%%%%%%%%%%%%%%%%%%%%%%%%%%%%%%%%%%%%%%%%%%%%%%%%%%%%%%%%
\subsection{Standard Error of a Contrast}
\label{subsec:standard-error-contrast}
%%%%%%%%%%%%%%%%%%%%%%%%%%%%%%%%%%%%%%%%%%%%%%%%%%%%%%%%%%%%

Under the equal-variance one-way ANOVA model,

\[
\boxed{
\operatorname{Var}(\widehat{L})
=
\sigma^2
\sum_{i=1}^{k}
\frac{
c_i^2
}{
n_i
}.
}
\]

Replacing \(\sigma^2\) by \(MS_E\),

\[
\boxed{
SE(\widehat{L})
=
\sqrt{
MS_E
\sum_{i=1}^{k}
\frac{
c_i^2
}{
n_i
}
}.
}
\]

Thus a \(t\)-statistic for

\[
H_0:L=0
\]

is

\[
\boxed{
T
=
\frac{
\widehat{L}
}{
SE(\widehat{L})
}.
}
\]

%%%%%%%%%%%%%%%%%%%%%%%%%%%%%%%%%%%%%%%%%%%%%%%%%%%%%%%%%%%%
\subsection{Orthogonal Contrasts}
\label{subsec:orthogonal-contrasts}
%%%%%%%%%%%%%%%%%%%%%%%%%%%%%%%%%%%%%%%%%%%%%%%%%%%%%%%%%%%%

In a balanced design, two contrasts with coefficient vectors

\[
(c_1,\ldots,c_k)
\]

and

\[
(d_1,\ldots,d_k)
\]

are orthogonal if

\[
\boxed{
\sum_{i=1}^{k}
c_id_i
=
0.
}
\]

For unequal group sizes, the appropriate orthogonality relation depends
on the contrast parameterization and weighting.

Orthogonal contrasts divide between-group variability into independent
components under the normal balanced model.

%%%%%%%%%%%%%%%%%%%%%%%%%%%%%%%%%%%%%%%%%%%%%%%%%%%%%%%%%%%%
\subsection{Number of Independent Contrasts}
\label{subsec:number-independent-contrasts}
%%%%%%%%%%%%%%%%%%%%%%%%%%%%%%%%%%%%%%%%%%%%%%%%%%%%%%%%%%%%

With \(k\) group means, the treatment space has

\[
\boxed{
k-1
}
\]

degrees of freedom.

Therefore at most

\[
\boxed{
k-1
}
\]

linearly independent contrasts can be formed.

A complete orthogonal set can partition the treatment sum of squares.

%%%%%%%%%%%%%%%%%%%%%%%%%%%%%%%%%%%%%%%%%%%%%%%%%%%%%%%%%%%%
\subsection{Omnibus Tests Versus Pairwise Comparisons}
\label{subsec:omnibus-versus-pairwise}
%%%%%%%%%%%%%%%%%%%%%%%%%%%%%%%%%%%%%%%%%%%%%%%%%%%%%%%%%%%%

The omnibus ANOVA test asks

\[
\boxed{
\text{Are all population means equal?}
}
\]

Pairwise comparisons ask questions such as

\[
\boxed{
\mu_i-\mu_j=0.
}
\]

A significant omnibus test does not automatically identify which pairs
differ.

Likewise, specific planned contrasts can be scientifically meaningful
even if they do not correspond to all possible pairwise comparisons.

%%%%%%%%%%%%%%%%%%%%%%%%%%%%%%%%%%%%%%%%%%%%%%%%%%%%%%%%%%%%
\subsection{Multiple Comparisons}
\label{subsec:anova-multiple-comparisons}
%%%%%%%%%%%%%%%%%%%%%%%%%%%%%%%%%%%%%%%%%%%%%%%%%%%%%%%%%%%%

With \(k\) groups, the number of pairwise comparisons is

\[
\boxed{
\binom{k}{2}
=
\frac{
k(k-1)
}{2}.
}
\]

Testing all of them individually at level

\[
\alpha
\]

inflates the probability of at least one false positive.

Thus multiplicity adjustment is often required.

%%%%%%%%%%%%%%%%%%%%%%%%%%%%%%%%%%%%%%%%%%%%%%%%%%%%%%%%%%%%
\subsection{Bonferroni Pairwise Comparisons}
\label{subsec:anova-bonferroni}
%%%%%%%%%%%%%%%%%%%%%%%%%%%%%%%%%%%%%%%%%%%%%%%%%%%%%%%%%%%%

Suppose \(m\) comparisons are planned and family-wise error rate should
be at most \(\alpha\).

Bonferroni uses individual significance level

\[
\boxed{
\alpha_{\text{individual}}
=
\frac{\alpha}{m}.
}
\]

Equivalently, individual confidence intervals may use level

\[
\boxed{
1-\frac{\alpha}{m}.
}
\]

Bonferroni is simple and valid under arbitrary dependence, though it may
be conservative.

%%%%%%%%%%%%%%%%%%%%%%%%%%%%%%%%%%%%%%%%%%%%%%%%%%%%%%%%%%%%
\subsection{Tukey's Honest Significant Difference}
\label{subsec:tukey-hsd}
%%%%%%%%%%%%%%%%%%%%%%%%%%%%%%%%%%%%%%%%%%%%%%%%%%%%%%%%%%%%

Tukey's HSD procedure is designed for simultaneous pairwise comparisons
among group means.

In a balanced one-way ANOVA, differences are compared using the
studentized range distribution.

A schematic threshold is

\[
\boxed{
|\overline{Y}_{i\cdot}
-
\overline{Y}_{j\cdot}|
>
q^*
\sqrt{
\frac{
MS_E
}{n}
}.
}
\]

Exact formulas depend on the studentized-range convention and whether
group sizes are equal.

%%%%%%%%%%%%%%%%%%%%%%%%%%%%%%%%%%%%%%%%%%%%%%%%%%%%%%%%%%%%
\subsection{Tukey--Kramer Procedure}
\label{subsec:tukey-kramer}
%%%%%%%%%%%%%%%%%%%%%%%%%%%%%%%%%%%%%%%%%%%%%%%%%%%%%%%%%%%%

For unequal group sizes, the Tukey--Kramer modification uses a pairwise
standard error involving

\[
\boxed{
\frac12
\left(
\frac1{n_i}
+
\frac1{n_j}
\right).
}
\]

This allows simultaneous pairwise comparisons while controlling
family-wise error under the classical model.

%%%%%%%%%%%%%%%%%%%%%%%%%%%%%%%%%%%%%%%%%%%%%%%%%%%%%%%%%%%%
\subsection{Planned Versus Post Hoc Comparisons}
\label{subsec:planned-versus-posthoc}
%%%%%%%%%%%%%%%%%%%%%%%%%%%%%%%%%%%%%%%%%%%%%%%%%%%%%%%%%%%%

Planned comparisons are specified before examining the data.

Post hoc comparisons are chosen after observing results.

Planned comparisons can focus statistical power on scientifically
motivated hypotheses.

Post hoc exploration requires greater attention to multiplicity because
the data influence which hypotheses are examined.

%%%%%%%%%%%%%%%%%%%%%%%%%%%%%%%%%%%%%%%%%%%%%%%%%%%%%%%%%%%%
\subsection{Effect Size in ANOVA}
\label{subsec:anova-effect-size}
%%%%%%%%%%%%%%%%%%%%%%%%%%%%%%%%%%%%%%%%%%%%%%%%%%%%%%%%%%%%

Statistical significance does not communicate how much variation is
associated with group membership.

An effect-size measure can supplement the \(F\)-test.

One common measure is eta-squared:

\[
\boxed{
\eta^2
=
\frac{
SS_A
}{
SS_T
}.
}
\]

This is the proportion of total sample variation associated with the
factor in the one-way decomposition.

%%%%%%%%%%%%%%%%%%%%%%%%%%%%%%%%%%%%%%%%%%%%%%%%%%%%%%%%%%%%
\subsection{Partial Eta-Squared}
\label{subsec:partial-eta-squared}
%%%%%%%%%%%%%%%%%%%%%%%%%%%%%%%%%%%%%%%%%%%%%%%%%%%%%%%%%%%%

In multifactor models, partial eta-squared for effect \(A\) is commonly

\[
\boxed{
\eta_p^2
=
\frac{
SS_A
}{
SS_A+SS_E
}.
}
\]

It measures the proportion of effect-plus-error variation attributable to
the specified effect.

Partial eta-squared and ordinary eta-squared are not interchangeable.

%%%%%%%%%%%%%%%%%%%%%%%%%%%%%%%%%%%%%%%%%%%%%%%%%%%%%%%%%%%%
\subsection{Omega-Squared Preview}
\label{subsec:omega-squared-preview}
%%%%%%%%%%%%%%%%%%%%%%%%%%%%%%%%%%%%%%%%%%%%%%%%%%%%%%%%%%%%

Eta-squared can be biased upward as an estimate of a population effect.

A common one-way omega-squared estimator is

\[
\boxed{
\omega^2
=
\frac{
SS_A
-
(k-1)MS_E
}{
SS_T+MS_E
}.
}
\]

Depending on the sample, the raw estimate may become slightly negative;
in descriptive reporting it is sometimes truncated at zero.

%%%%%%%%%%%%%%%%%%%%%%%%%%%%%%%%%%%%%%%%%%%%%%%%%%%%%%%%%%%%
\subsection{Welch's ANOVA}
\label{subsec:welch-anova}
%%%%%%%%%%%%%%%%%%%%%%%%%%%%%%%%%%%%%%%%%%%%%%%%%%%%%%%%%%%%

If population variances differ substantially, ordinary one-way ANOVA may
not be appropriate.

Welch's ANOVA allows unequal variances and does not pool them into one
common \(MS_E\).

Groups are weighted according to their estimated variances.

The resulting statistic uses an approximate \(F\)-distribution with
adjusted degrees of freedom.

%%%%%%%%%%%%%%%%%%%%%%%%%%%%%%%%%%%%%%%%%%%%%%%%%%%%%%%%%%%%
\subsection{When Welch's ANOVA Is Useful}
\label{subsec:when-welch-anova}
%%%%%%%%%%%%%%%%%%%%%%%%%%%%%%%%%%%%%%%%%%%%%%%%%%%%%%%%%%%%

Welch's ANOVA is especially useful when:

\[
\boxed{
\text{group variances differ},
}
\]

and particularly when

\[
\boxed{
\text{group sizes are also unequal}.
}
\]

This combination can seriously distort the classical equal-variance
ANOVA.

%%%%%%%%%%%%%%%%%%%%%%%%%%%%%%%%%%%%%%%%%%%%%%%%%%%%%%%%%%%%
\subsection{Testing Equality of Variances}
\label{subsec:testing-equality-variances-anova}
%%%%%%%%%%%%%%%%%%%%%%%%%%%%%%%%%%%%%%%%%%%%%%%%%%%%%%%%%%%%

Formal tests of equal variances exist, including:

\[
\boxed{
\text{Bartlett's test},
}
\]

\[
\boxed{
\text{Levene's test},
}
\]

and

\[
\boxed{
\text{Brown--Forsythe-type tests}.
}
\]

Bartlett's test is sensitive to nonnormality.

Levene-type procedures are generally more robust.

However, choosing an ANOVA method solely from a preliminary variance test
can create additional inferential complications.

%%%%%%%%%%%%%%%%%%%%%%%%%%%%%%%%%%%%%%%%%%%%%%%%%%%%%%%%%%%%
\subsection{Kruskal--Wallis Test Preview}
\label{subsec:kruskal-wallis-preview}
%%%%%%%%%%%%%%%%%%%%%%%%%%%%%%%%%%%%%%%%%%%%%%%%%%%%%%%%%%%%

The Kruskal--Wallis test is a rank-based alternative for comparing
multiple independent groups.

It does not require normality.

The observations are replaced by ranks, and the groups are compared
through rank sums.

Under appropriate conditions, the test statistic has an approximate
chi-square distribution.

%%%%%%%%%%%%%%%%%%%%%%%%%%%%%%%%%%%%%%%%%%%%%%%%%%%%%%%%%%%%
\subsection{Interpreting Kruskal--Wallis Carefully}
\label{subsec:kruskal-wallis-interpretation}
%%%%%%%%%%%%%%%%%%%%%%%%%%%%%%%%%%%%%%%%%%%%%%%%%%%%%%%%%%%%

The Kruskal--Wallis test is not automatically a test of equal medians.

Its precise interpretation depends on assumptions about the shapes of the
group distributions.

Under similarly shaped distributions differing mainly in location, it can
be interpreted as a test of location differences.

%%%%%%%%%%%%%%%%%%%%%%%%%%%%%%%%%%%%%%%%%%%%%%%%%%%%%%%%%%%%
\subsection{Two-Way ANOVA}
\label{subsec:two-way-anova}
%%%%%%%%%%%%%%%%%%%%%%%%%%%%%%%%%%%%%%%%%%%%%%%%%%%%%%%%%%%%

Two-way ANOVA studies two categorical factors simultaneously.

Suppose factor \(A\) has levels

\[
i=1,\ldots,a,
\]

and factor \(B\) has levels

\[
j=1,\ldots,b.
\]

With replicate \(k\), write

\[
Y_{ijk}.
\]

A standard model is

\[
\boxed{
Y_{ijk}
=
\mu
+
\alpha_i
+
\beta_j
+
(\alpha\beta)_{ij}
+
\varepsilon_{ijk}.
}
\]

%%%%%%%%%%%%%%%%%%%%%%%%%%%%%%%%%%%%%%%%%%%%%%%%%%%%%%%%%%%%
\subsection{Main Effect of Factor \(A\)}
\label{subsec:main-effect-a}
%%%%%%%%%%%%%%%%%%%%%%%%%%%%%%%%%%%%%%%%%%%%%%%%%%%%%%%%%%%%

The main effect of \(A\) describes how the response changes across levels
of \(A\) after averaging over levels of \(B\).

The null hypothesis is

\[
\boxed{
H_0^{(A)}:
\alpha_1=\cdots=\alpha_a=0
}
\]

under an appropriate sum-to-zero parameterization.

Conceptually,

\[
\boxed{
\text{no main effect of }A
}
\]

means marginal means do not differ across levels of \(A\).

%%%%%%%%%%%%%%%%%%%%%%%%%%%%%%%%%%%%%%%%%%%%%%%%%%%%%%%%%%%%
\subsection{Main Effect of Factor \(B\)}
\label{subsec:main-effect-b}
%%%%%%%%%%%%%%%%%%%%%%%%%%%%%%%%%%%%%%%%%%%%%%%%%%%%%%%%%%%%

Similarly, the main effect of \(B\) compares marginal means across levels
of factor \(B\).

The null hypothesis is

\[
\boxed{
H_0^{(B)}:
\beta_1=\cdots=\beta_b=0.
}
\]

%%%%%%%%%%%%%%%%%%%%%%%%%%%%%%%%%%%%%%%%%%%%%%%%%%%%%%%%%%%%
\subsection{Interaction Effect}
\label{subsec:anova-interaction}
%%%%%%%%%%%%%%%%%%%%%%%%%%%%%%%%%%%%%%%%%%%%%%%%%%%%%%%%%%%%

Interaction occurs when the effect of one factor depends on the level of
the other factor.

In the model,

\[
\boxed{
(\alpha\beta)_{ij}
}
\]

represents the interaction between level \(i\) of \(A\) and level \(j\)
of \(B\).

The interaction null is

\[
\boxed{
H_0^{(AB)}:
(\alpha\beta)_{ij}=0
\text{ for all }i,j.
}
\]

%%%%%%%%%%%%%%%%%%%%%%%%%%%%%%%%%%%%%%%%%%%%%%%%%%%%%%%%%%%%
\subsection{Worked Example: Interaction}
\label{subsec:worked-anova-interaction}
%%%%%%%%%%%%%%%%%%%%%%%%%%%%%%%%%%%%%%%%%%%%%%%%%%%%%%%%%%%%

\begin{workedexamplebox}{Treatment Effect Depends on Group}

Suppose two treatments are tested in two populations.

The mean responses are

\[
\begin{array}{c|cc}
&
\text{Treatment 1}
&
\text{Treatment 2}
\\
\hline
\text{Population A}
&
10
&
20
\\
\text{Population B}
&
30
&
25
\end{array}
\]

For population A, treatment \(2\) increases the response by

\[
10.
\]

For population B, treatment \(2\) decreases the response by

\[
5.
\]

\begin{answerbox}
\[
\boxed{
\text{The treatment effect depends on population, indicating
interaction.}
}
\]
\end{answerbox}

\end{workedexamplebox}

%%%%%%%%%%%%%%%%%%%%%%%%%%%%%%%%%%%%%%%%%%%%%%%%%%%%%%%%%%%%
\subsection{Interaction Plots}
\label{subsec:interaction-plots}
%%%%%%%%%%%%%%%%%%%%%%%%%%%%%%%%%%%%%%%%%%%%%%%%%%%%%%%%%%%%

An interaction plot graphs mean response across levels of one factor,
with separate lines for levels of the second factor.

Roughly parallel lines suggest little interaction.

Strongly nonparallel or crossing lines suggest interaction.

The plot is descriptive; formal inference still uses the appropriate
model.

%%%%%%%%%%%%%%%%%%%%%%%%%%%%%%%%%%%%%%%%%%%%%%%%%%%%%%%%%%%%
\subsection{Why Interaction Changes Main-Effect Interpretation}
\label{subsec:interaction-main-effects}
%%%%%%%%%%%%%%%%%%%%%%%%%%%%%%%%%%%%%%%%%%%%%%%%%%%%%%%%%%%%

If interaction is strong, a single averaged main effect may hide
important differences.

For example, one treatment may help one subgroup while harming another.

Thus when interaction is present, interpretation often focuses on:

\[
\boxed{
\text{simple effects}
}
\]

or

\[
\boxed{
\text{factor effects at specific levels of the other factor}.
}
\]

%%%%%%%%%%%%%%%%%%%%%%%%%%%%%%%%%%%%%%%%%%%%%%%%%%%%%%%%%%%%
\subsection{Balanced Two-Way Designs}
\label{subsec:balanced-two-way-anova}
%%%%%%%%%%%%%%%%%%%%%%%%%%%%%%%%%%%%%%%%%%%%%%%%%%%%%%%%%%%%

A two-way design is balanced if every factor combination has the same
number of observations.

If there are

\[
n
\]

replicates per cell, then the total sample size is

\[
\boxed{
N=abn.
}
\]

Balanced designs simplify sums-of-squares decompositions and produce
orthogonality among many effects.

%%%%%%%%%%%%%%%%%%%%%%%%%%%%%%%%%%%%%%%%%%%%%%%%%%%%%%%%%%%%
\subsection{Two-Way ANOVA Degrees of Freedom}
\label{subsec:two-way-anova-df}
%%%%%%%%%%%%%%%%%%%%%%%%%%%%%%%%%%%%%%%%%%%%%%%%%%%%%%%%%%%%

For a balanced two-way factorial design with replication:

\[
\boxed{
df_A
=
a-1,
}
\]

\[
\boxed{
df_B
=
b-1,
}
\]

\[
\boxed{
df_{AB}
=
(a-1)(b-1),
}
\]

and

\[
\boxed{
df_E
=
ab(n-1).
}
\]

The total degrees of freedom are

\[
\boxed{
abn-1.
}
\]

%%%%%%%%%%%%%%%%%%%%%%%%%%%%%%%%%%%%%%%%%%%%%%%%%%%%%%%%%%%%
\subsection{Two-Way ANOVA Table}
\label{subsec:two-way-anova-table}
%%%%%%%%%%%%%%%%%%%%%%%%%%%%%%%%%%%%%%%%%%%%%%%%%%%%%%%%%%%%

A balanced two-way ANOVA table has the structure

\[
\begin{array}{c|c|c|c|c}
\text{Source}
&
\text{df}
&
\text{SS}
&
\text{MS}
&
F
\\
\hline
A
&
a-1
&
SS_A
&
MS_A
&
MS_A/MS_E
\\
B
&
b-1
&
SS_B
&
MS_B
&
MS_B/MS_E
\\
A\times B
&
(a-1)(b-1)
&
SS_{AB}
&
MS_{AB}
&
MS_{AB}/MS_E
\\
\text{Error}
&
ab(n-1)
&
SS_E
&
MS_E
&
\\
\text{Total}
&
abn-1
&
SS_T
&
&
\end{array}
\]

Each \(F\)-ratio compares one modeled source of variation with residual
variation.

%%%%%%%%%%%%%%%%%%%%%%%%%%%%%%%%%%%%%%%%%%%%%%%%%%%%%%%%%%%%
\subsection{Factorial Designs}
\label{subsec:factorial-designs-anova}
%%%%%%%%%%%%%%%%%%%%%%%%%%%%%%%%%%%%%%%%%%%%%%%%%%%%%%%%%%%%

A factorial design studies all combinations of levels of two or more
factors.

For example, a

\[
2\times3
\]

factorial design has:

\[
\boxed{
2
\text{ levels of factor }A
}
\]

and

\[
\boxed{
3
\text{ levels of factor }B.
}
\]

Therefore there are

\[
\boxed{
6
}
\]

treatment combinations.

%%%%%%%%%%%%%%%%%%%%%%%%%%%%%%%%%%%%%%%%%%%%%%%%%%%%%%%%%%%%
\subsection{Advantages of Factorial Designs}
\label{subsec:factorial-design-advantages}
%%%%%%%%%%%%%%%%%%%%%%%%%%%%%%%%%%%%%%%%%%%%%%%%%%%%%%%%%%%%

Factorial designs allow researchers to estimate:

\[
\boxed{
\text{main effects}
}
\]

and

\[
\boxed{
\text{interactions}
}
\]

within one experiment.

They can be much more informative than studying factors separately.

This is one reason factorial experimentation is central to experimental
design.

%%%%%%%%%%%%%%%%%%%%%%%%%%%%%%%%%%%%%%%%%%%%%%%%%%%%%%%%%%%%
\subsection{Blocking and ANOVA}
\label{subsec:blocking-anova}
%%%%%%%%%%%%%%%%%%%%%%%%%%%%%%%%%%%%%%%%%%%%%%%%%%%%%%%%%%%%

Blocking incorporates a known nuisance source of variability.

A randomized complete block model may be written

\[
\boxed{
Y_{ij}
=
\mu
+
\tau_i
+
\beta_j
+
\varepsilon_{ij},
}
\]

where

\[
\tau_i
\]

is the treatment effect and

\[
\beta_j
\]

is the block effect.

Blocking can reduce the residual variance used to compare treatments.

%%%%%%%%%%%%%%%%%%%%%%%%%%%%%%%%%%%%%%%%%%%%%%%%%%%%%%%%%%%%
\subsection{Randomized Complete Block Design}
\label{subsec:randomized-complete-block-design}
%%%%%%%%%%%%%%%%%%%%%%%%%%%%%%%%%%%%%%%%%%%%%%%%%%%%%%%%%%%%

In a randomized complete block design:

\begin{enumerate}
    \item experimental units are grouped into homogeneous blocks;
    \item every treatment appears once or equally often in each block;
    \item treatment assignments are randomized within blocks.
\end{enumerate}

This allows treatment comparisons after accounting for block-to-block
variation.

%%%%%%%%%%%%%%%%%%%%%%%%%%%%%%%%%%%%%%%%%%%%%%%%%%%%%%%%%%%%
\subsection{Repeated-Measures ANOVA Preview}
\label{subsec:repeated-measures-anova}
%%%%%%%%%%%%%%%%%%%%%%%%%%%%%%%%%%%%%%%%%%%%%%%%%%%%%%%%%%%%

Repeated-measures designs observe the same experimental unit under
multiple conditions or times.

Measurements from the same unit are correlated.

Therefore ordinary independent-groups ANOVA is inappropriate.

Repeated-measures ANOVA separates:

\[
\boxed{
\text{between-subject variability}
}
\]

from

\[
\boxed{
\text{within-subject treatment variability}.
}
\]

%%%%%%%%%%%%%%%%%%%%%%%%%%%%%%%%%%%%%%%%%%%%%%%%%%%%%%%%%%%%
\subsection{Sphericity Preview}
\label{subsec:sphericity-preview}
%%%%%%%%%%%%%%%%%%%%%%%%%%%%%%%%%%%%%%%%%%%%%%%%%%%%%%%%%%%%

Classical repeated-measures ANOVA may require a covariance condition
called sphericity.

Informally, sphericity requires variances of pairwise condition
differences to be equal.

Violations affect the reference \(F\)-distribution.

Corrections such as Greenhouse--Geisser or Huynh--Feldt are commonly
used.

%%%%%%%%%%%%%%%%%%%%%%%%%%%%%%%%%%%%%%%%%%%%%%%%%%%%%%%%%%%%
\subsection{Mixed Models Preview}
\label{subsec:mixed-models-anova-preview}
%%%%%%%%%%%%%%%%%%%%%%%%%%%%%%%%%%%%%%%%%%%%%%%%%%%%%%%%%%%%

Repeated, clustered, or hierarchical data are often handled more
flexibly with mixed-effects models.

A mixed model contains both:

\[
\boxed{
\text{fixed effects}
}
\]

and

\[
\boxed{
\text{random effects}.
}
\]

Mixed models generalize classical ANOVA and regression to correlated and
multilevel data.

%%%%%%%%%%%%%%%%%%%%%%%%%%%%%%%%%%%%%%%%%%%%%%%%%%%%%%%%%%%%
\subsection{Fixed Effects}
\label{subsec:anova-fixed-effects}
%%%%%%%%%%%%%%%%%%%%%%%%%%%%%%%%%%%%%%%%%%%%%%%%%%%%%%%%%%%%

A factor is treated as fixed when the levels themselves are the specific
levels of interest.

For example, if three particular instructional methods are being
compared, those methods may be treated as fixed effects.

Inference concerns differences among those particular levels.

%%%%%%%%%%%%%%%%%%%%%%%%%%%%%%%%%%%%%%%%%%%%%%%%%%%%%%%%%%%%
\subsection{Random Effects}
\label{subsec:anova-random-effects}
%%%%%%%%%%%%%%%%%%%%%%%%%%%%%%%%%%%%%%%%%%%%%%%%%%%%%%%%%%%%

A factor may be treated as random when its observed levels are regarded
as a sample from a larger population of possible levels.

A simple random-effects model is

\[
\boxed{
Y_{ij}
=
\mu
+
A_i
+
\varepsilon_{ij},
}
\]

where

\[
A_i
\]

is itself random.

One may assume

\[
\boxed{
A_i
\sim
N(0,\sigma_A^2).
}
\]

The parameter of interest may then be a variance component rather than
specific level means.

%%%%%%%%%%%%%%%%%%%%%%%%%%%%%%%%%%%%%%%%%%%%%%%%%%%%%%%%%%%%
\subsection{Variance Components}
\label{subsec:variance-components}
%%%%%%%%%%%%%%%%%%%%%%%%%%%%%%%%%%%%%%%%%%%%%%%%%%%%%%%%%%%%

In a random-effects model,

\[
\operatorname{Var}(Y_{ij})
\]

may decompose into several sources.

For example,

\[
\boxed{
\operatorname{Var}(Y_{ij})
=
\sigma_A^2
+
\sigma^2.
}
\]

Here,

\[
\sigma_A^2
\]

represents between-group variance and

\[
\sigma^2
\]

represents within-group variance.

%%%%%%%%%%%%%%%%%%%%%%%%%%%%%%%%%%%%%%%%%%%%%%%%%%%%%%%%%%%%
\subsection{Intraclass Correlation Preview}
\label{subsec:intraclass-correlation-anova}
%%%%%%%%%%%%%%%%%%%%%%%%%%%%%%%%%%%%%%%%%%%%%%%%%%%%%%%%%%%%

For a simple random-intercept model, observations in the same group share
the random effect \(A_i\).

The intraclass correlation coefficient is

\[
\boxed{
\rho
=
\frac{
\sigma_A^2
}{
\sigma_A^2+\sigma^2
}.
}
\]

This measures similarity among observations from the same cluster.

%%%%%%%%%%%%%%%%%%%%%%%%%%%%%%%%%%%%%%%%%%%%%%%%%%%%%%%%%%%%
\subsection{Nested Factors}
\label{subsec:nested-factors-anova}
%%%%%%%%%%%%%%%%%%%%%%%%%%%%%%%%%%%%%%%%%%%%%%%%%%%%%%%%%%%%

Factor \(B\) is nested within factor \(A\) when each level of \(B\)
occurs within only one level of \(A\).

For example,

\[
\boxed{
\text{students nested within classrooms}
}
\]

or

\[
\boxed{
\text{machines nested within factories}.
}
\]

Nested designs differ structurally from crossed factorial designs.

%%%%%%%%%%%%%%%%%%%%%%%%%%%%%%%%%%%%%%%%%%%%%%%%%%%%%%%%%%%%
\subsection{Crossed Versus Nested Factors}
\label{subsec:crossed-versus-nested}
%%%%%%%%%%%%%%%%%%%%%%%%%%%%%%%%%%%%%%%%%%%%%%%%%%%%%%%%%%%%

Factors are crossed when every level of one factor can occur with every
level of another factor.

For example,

\[
\boxed{
\text{treatment}
\times
\text{sex}
}
\]

may be crossed.

Factors are nested when levels belong uniquely inside another factor.

Correct identification affects the ANOVA model and appropriate error
terms.

%%%%%%%%%%%%%%%%%%%%%%%%%%%%%%%%%%%%%%%%%%%%%%%%%%%%%%%%%%%%
\subsection{ANCOVA Preview}
\label{subsec:ancova-preview}
%%%%%%%%%%%%%%%%%%%%%%%%%%%%%%%%%%%%%%%%%%%%%%%%%%%%%%%%%%%%

Analysis of covariance, abbreviated ANCOVA, combines categorical
treatment factors with quantitative covariates.

A simple model is

\[
\boxed{
Y
=
\beta_0
+
\beta_1X
+
\tau_i
+
\varepsilon,
}
\]

where

\[
X
\]

is a quantitative covariate and

\[
\tau_i
\]

represents group effects.

ANCOVA is therefore a direct extension of multiple regression.

%%%%%%%%%%%%%%%%%%%%%%%%%%%%%%%%%%%%%%%%%%%%%%%%%%%%%%%%%%%%
\subsection{Why Use ANCOVA?}
\label{subsec:why-use-ancova}
%%%%%%%%%%%%%%%%%%%%%%%%%%%%%%%%%%%%%%%%%%%%%%%%%%%%%%%%%%%%

A covariate may explain response variation unrelated to treatment.

Adjusting for it can:

\[
\boxed{
\text{reduce residual variability},
}
\]

\[
\boxed{
\text{increase precision},
}
\]

and

\[
\boxed{
\text{adjust comparisons for baseline differences}.
}
\]

Causal interpretation still depends on study design and appropriate
covariate selection.

%%%%%%%%%%%%%%%%%%%%%%%%%%%%%%%%%%%%%%%%%%%%%%%%%%%%%%%%%%%%
\subsection{Equal-Slopes Assumption in Classical ANCOVA}
\label{subsec:ancova-equal-slopes}
%%%%%%%%%%%%%%%%%%%%%%%%%%%%%%%%%%%%%%%%%%%%%%%%%%%%%%%%%%%%

A simple ANCOVA model without interaction assumes the covariate slope is
the same across groups.

That is,

\[
\boxed{
\text{no group}\times\text{covariate interaction}.
}
\]

If slopes differ substantially, the model should include an interaction
term rather than forcing parallel regression lines.

%%%%%%%%%%%%%%%%%%%%%%%%%%%%%%%%%%%%%%%%%%%%%%%%%%%%%%%%%%%%
\subsection{Unbalanced ANOVA}
\label{subsec:unbalanced-anova}
%%%%%%%%%%%%%%%%%%%%%%%%%%%%%%%%%%%%%%%%%%%%%%%%%%%%%%%%%%%%

A design is unbalanced when group or cell sizes differ.

In one-way ANOVA, the sums-of-squares decomposition remains relatively
straightforward.

In multifactor ANOVA, unbalanced data introduce complications because
effects may not be orthogonal.

Different sums-of-squares definitions can then produce different
numerical results.

%%%%%%%%%%%%%%%%%%%%%%%%%%%%%%%%%%%%%%%%%%%%%%%%%%%%%%%%%%%%
\subsection{Type I, II, and III Sums of Squares Preview}
\label{subsec:type-sums-of-squares}
%%%%%%%%%%%%%%%%%%%%%%%%%%%%%%%%%%%%%%%%%%%%%%%%%%%%%%%%%%%%

In unbalanced factorial models, statistical software may report
different sums of squares.

\textbf{Type I sums of squares} are sequential and depend on the order
terms enter the model.

\textbf{Type II sums of squares} test main effects after other main
effects but typically not interactions containing that effect.

\textbf{Type III sums of squares} test each effect conditional on all
other terms in the specified model under a chosen parameterization.

These methods answer different questions and should not be selected
mechanically.

%%%%%%%%%%%%%%%%%%%%%%%%%%%%%%%%%%%%%%%%%%%%%%%%%%%%%%%%%%%%
\subsection{ANOVA and the General Linear Model}
\label{subsec:anova-general-linear-model}
%%%%%%%%%%%%%%%%%%%%%%%%%%%%%%%%%%%%%%%%%%%%%%%%%%%%%%%%%%%%

Regression and ANOVA can both be written as

\[
\boxed{
\mathbf{y}
=
X\boldsymbol{\beta}
+
\boldsymbol{\varepsilon}.
}
\]

Continuous predictors, categorical indicators, interactions, and
polynomial terms all enter through the design matrix \(X\).

Therefore the distinction between ANOVA and regression is often mainly
one of interpretation and predictor type rather than underlying
mathematics.

%%%%%%%%%%%%%%%%%%%%%%%%%%%%%%%%%%%%%%%%%%%%%%%%%%%%%%%%%%%%
\subsection{General Linear Hypotheses}
\label{subsec:anova-general-linear-hypotheses}
%%%%%%%%%%%%%%%%%%%%%%%%%%%%%%%%%%%%%%%%%%%%%%%%%%%%%%%%%%%%

Many ANOVA hypotheses can be written

\[
\boxed{
H_0:
R\boldsymbol{\beta}
=
\mathbf{r}.
}
\]

For example, equality of several group means corresponds to multiple
linear restrictions.

This formulation unifies:

\[
\boxed{
\text{ANOVA},
}
\]

\[
\boxed{
\text{regression coefficient tests},
}
\]

and

\[
\boxed{
\text{contrasts}.
}
\]

%%%%%%%%%%%%%%%%%%%%%%%%%%%%%%%%%%%%%%%%%%%%%%%%%%%%%%%%%%%%
\subsection{Partial \(F\)-Tests}
\label{subsec:partial-f-tests-anova}
%%%%%%%%%%%%%%%%%%%%%%%%%%%%%%%%%%%%%%%%%%%%%%%%%%%%%%%%%%%%

Suppose a reduced model is nested inside a full model.

Let

\[
SSE_R
\]

be the reduced-model error sum of squares and

\[
SSE_F
\]

the full-model error sum of squares.

If the full model adds \(q\) parameters, then

\[
\boxed{
F
=
\frac{
(SSE_R-SSE_F)/q
}{
SSE_F/df_F
}.
}
\]

Under the classical null,

\[
F
\]

has an \(F\)-distribution with numerator degrees of freedom \(q\) and
denominator degrees of freedom \(df_F\).

%%%%%%%%%%%%%%%%%%%%%%%%%%%%%%%%%%%%%%%%%%%%%%%%%%%%%%%%%%%%
\subsection{Connection Between Partial \(F\)-Tests and ANOVA}
\label{subsec:partial-f-anova-connection}
%%%%%%%%%%%%%%%%%%%%%%%%%%%%%%%%%%%%%%%%%%%%%%%%%%%%%%%%%%%%

A factor in ANOVA may correspond to several regression coefficients.

Testing that factor means testing those coefficients jointly.

Thus a factor with \(a\) levels usually contributes

\[
a-1
\]

degrees of freedom.

The corresponding ANOVA \(F\)-test is a partial \(F\)-test for the group
of indicator coefficients.

%%%%%%%%%%%%%%%%%%%%%%%%%%%%%%%%%%%%%%%%%%%%%%%%%%%%%%%%%%%%
\subsection{Effect Size Versus Significance}
\label{subsec:anova-effect-size-significance}
%%%%%%%%%%%%%%%%%%%%%%%%%%%%%%%%%%%%%%%%%%%%%%%%%%%%%%%%%%%%

A significant \(F\)-test does not indicate whether the effect is large.

With very large \(N\), very small differences can become significant.

Therefore ANOVA results should ideally include:

\[
\boxed{
\text{group estimates},
}
\]

\[
\boxed{
\text{confidence intervals},
}
\]

\[
\boxed{
\text{effect sizes},
}
\]

and

\[
\boxed{
\text{scientific interpretation}.
}
\]

%%%%%%%%%%%%%%%%%%%%%%%%%%%%%%%%%%%%%%%%%%%%%%%%%%%%%%%%%%%%
\subsection{Power in One-Way ANOVA Preview}
\label{subsec:anova-power-preview}
%%%%%%%%%%%%%%%%%%%%%%%%%%%%%%%%%%%%%%%%%%%%%%%%%%%%%%%%%%%%

ANOVA power depends on:

\[
\boxed{
\text{number of groups},
}
\]

\[
\boxed{
\text{sample sizes},
}
\]

\[
\boxed{
\text{within-group variance},
}
\]

and

\[
\boxed{
\text{separation among population means}.
}
\]

Under alternatives, the \(F\)-statistic is related to a noncentral
\(F\)-distribution.

This provides the basis for formal ANOVA power calculations.

%%%%%%%%%%%%%%%%%%%%%%%%%%%%%%%%%%%%%%%%%%%%%%%%%%%%%%%%%%%%
\subsection{Cohen's \(f\) Preview}
\label{subsec:cohens-f-anova}
%%%%%%%%%%%%%%%%%%%%%%%%%%%%%%%%%%%%%%%%%%%%%%%%%%%%%%%%%%%%

A standardized one-way ANOVA effect size may be written

\[
\boxed{
f
=
\sqrt{
\frac{
\eta^2
}{
1-\eta^2
}
}.
}
\]

This quantity is useful in some power and sample-size formulas.

Interpretation should depend on the application rather than rigid
universal categories.

%%%%%%%%%%%%%%%%%%%%%%%%%%%%%%%%%%%%%%%%%%%%%%%%%%%%%%%%%%%%
\subsection{ANOVA Modeling Workflow}
\label{subsec:anova-workflow}
%%%%%%%%%%%%%%%%%%%%%%%%%%%%%%%%%%%%%%%%%%%%%%%%%%%%%%%%%%%%

A practical ANOVA workflow is:

\begin{enumerate}
    \item identify the response variable;
    \item identify the categorical factor or factors;
    \item identify the experimental or observational design;
    \item inspect group sample sizes;
    \item calculate group means and spreads;
    \item visualize the group distributions;
    \item specify the ANOVA model;
    \item compute or fit the model;
    \item examine the omnibus \(F\)-test;
    \item inspect residual diagnostics;
    \item evaluate variance assumptions;
    \item examine interactions before interpreting main effects;
    \item perform planned contrasts or adjusted post hoc comparisons;
    \item report confidence intervals and effect sizes;
    \item interpret results in the context of the study design.
\end{enumerate}

%%%%%%%%%%%%%%%%%%%%%%%%%%%%%%%%%%%%%%%%%%%%%%%%%%%%%%%%%%%%
\subsection{Common Errors}
\label{subsec:anova-common-errors}
%%%%%%%%%%%%%%%%%%%%%%%%%%%%%%%%%%%%%%%%%%%%%%%%%%%%%%%%%%%%

\subsubsection{Performing Many Independent \(t\)-Tests Instead of ANOVA}

With many groups, repeated unadjusted pairwise tests inflate Type I error.

An omnibus ANOVA followed by appropriate comparisons provides a more
structured approach.

%%%%%%%%%%%%%%%%%%%%%%%%%%%%%%%%%%%%%%%%%%%%%%%%%%%%%%%%%%%%
\subsubsection{Interpreting a Significant ANOVA as Meaning Every Group Differs}

A significant \(F\)-test means only

\[
\boxed{
\text{not all means are equal}.
}
\]

Some group means may still be equal or nearly equal.

%%%%%%%%%%%%%%%%%%%%%%%%%%%%%%%%%%%%%%%%%%%%%%%%%%%%%%%%%%%%
\subsubsection{Interpreting a Nonsignificant ANOVA as Proof of Equality}

Failure to reject may reflect low power.

Equality requires a different scientific claim and possibly equivalence
methods.

%%%%%%%%%%%%%%%%%%%%%%%%%%%%%%%%%%%%%%%%%%%%%%%%%%%%%%%%%%%%
\subsubsection{Ignoring Unequal Variances}

Strong heteroskedasticity, especially with unequal group sizes, can make
the classical \(F\)-test unreliable.

Welch's ANOVA or another method may be more appropriate.

%%%%%%%%%%%%%%%%%%%%%%%%%%%%%%%%%%%%%%%%%%%%%%%%%%%%%%%%%%%%
\subsubsection{Ignoring Independence}

ANOVA does not fix repeated-measures or clustered dependence.

The model must reflect the data structure.

%%%%%%%%%%%%%%%%%%%%%%%%%%%%%%%%%%%%%%%%%%%%%%%%%%%%%%%%%%%%
\subsubsection{Testing Variance Equality and Then Automatically Choosing a Procedure}

A preliminary variance test changes the overall analysis strategy and can
have poor power.

Choice of method should consider study design, variance patterns, sample
sizes, and robustness rather than only one preliminary \(p\)-value.

%%%%%%%%%%%%%%%%%%%%%%%%%%%%%%%%%%%%%%%%%%%%%%%%%%%%%%%%%%%%
\subsubsection{Ignoring Interaction}

In factorial ANOVA, strong interaction can make averaged main effects
misleading.

Interaction should usually be examined before simple main-effect
interpretation.

%%%%%%%%%%%%%%%%%%%%%%%%%%%%%%%%%%%%%%%%%%%%%%%%%%%%%%%%%%%%
\subsubsection{Treating Group Labels as Numeric Predictors}

If groups are nominal categories, coding them as

\[
1,2,3
\]

and fitting one linear slope imposes an artificial quantitative ordering.

Indicator variables are generally appropriate.

%%%%%%%%%%%%%%%%%%%%%%%%%%%%%%%%%%%%%%%%%%%%%%%%%%%%%%%%%%%%
\subsubsection{Interpreting ANOVA as Causal Automatically}

A significant group effect in observational data does not prove that
group membership causes the response difference.

Causal interpretation requires appropriate design and assumptions.

%%%%%%%%%%%%%%%%%%%%%%%%%%%%%%%%%%%%%%%%%%%%%%%%%%%%%%%%%%%%
\subsubsection{Ignoring Multiplicity in Post Hoc Analysis}

Selecting and testing many comparisons after seeing the data can
substantially inflate false positive rates.

Use simultaneous or multiplicity-aware procedures.

%%%%%%%%%%%%%%%%%%%%%%%%%%%%%%%%%%%%%%%%%%%%%%%%%%%%%%%%%%%%
\subsubsection{Deleting Outliers Automatically}

An extreme observation may indicate:

\[
\boxed{
\text{data error},
}
\]

\[
\boxed{
\text{model failure},
}
\]

or

\[
\boxed{
\text{genuine scientific variation}.
}
\]

Investigate before removal.

%%%%%%%%%%%%%%%%%%%%%%%%%%%%%%%%%%%%%%%%%%%%%%%%%%%%%%%%%%%%
\subsubsection{Reporting Only the ANOVA \(p\)-Value}

An omnibus \(p\)-value does not communicate:

\[
\boxed{
\text{which groups differ},
}
\]

\[
\boxed{
\text{effect magnitude},
}
\]

or

\[
\boxed{
\text{uncertainty}.
}
\]

Report group estimates and relevant effect sizes.

%%%%%%%%%%%%%%%%%%%%%%%%%%%%%%%%%%%%%%%%%%%%%%%%%%%%%%%%%%%%
\subsection{Applications}
\label{subsec:anova-applications}
%%%%%%%%%%%%%%%%%%%%%%%%%%%%%%%%%%%%%%%%%%%%%%%%%%%%%%%%%%%%

\begin{applicationbox}

\textbf{Experimental Science.}
ANOVA compares multiple treatments while partitioning experimental
variation into interpretable sources.

\medskip

\textbf{Agriculture.}
Treatments, blocks, fields, fertilizers, irrigation levels, and crop
varieties are naturally studied through ANOVA designs.

\medskip

\textbf{Engineering.}
Factorial ANOVA evaluates how operating conditions and component choices
affect system performance.

\medskip

\textbf{Education Research.}
ANOVA compares instructional methods, programs, grade levels, and other
grouped educational outcomes.

\medskip

\textbf{Environmental Science.}
Measurements may be compared across sites, seasons, depths, treatments,
or geographic regions.

\medskip

\textbf{Medicine.}
ANOVA compares multiple treatments, doses, groups, or repeated
measurement conditions.

\medskip

\textbf{Psychology and Social Science.}
Factorial and repeated-measures ANOVA are widely used for experimental
and behavioral designs.

\medskip

\textbf{Industrial Experimentation.}
ANOVA forms the inferential foundation for factorial design, process
improvement, and design of experiments.

\end{applicationbox}

%%%%%%%%%%%%%%%%%%%%%%%%%%%%%%%%%%%%%%%%%%%%%%%%%%%%%%%%%%%%
\subsection{Exercises}
\label{subsec:anova-exercises}
%%%%%%%%%%%%%%%%%%%%%%%%%%%%%%%%%%%%%%%%%%%%%%%%%%%%%%%%%%%%

\begin{exercise}[1]
State the null hypothesis for one-way ANOVA with \(k\) groups.
\end{exercise}

\begin{exercise}[1]
Define the grand mean.
\end{exercise}

\begin{exercise}[1]
Define the total sum of squares.
\end{exercise}

\begin{exercise}[1]
Define the between-group sum of squares.
\end{exercise}

\begin{exercise}[1]
Define the within-group sum of squares.
\end{exercise}

\begin{exercise}[1]
State the ANOVA sum-of-squares decomposition.
\end{exercise}

\begin{exercise}[1]
State the degrees of freedom for a one-way ANOVA.
\end{exercise}

\begin{exercise}[1]
Define the ANOVA \(F\)-statistic.
\end{exercise}

\begin{exercise}[1]
State the classical one-way ANOVA assumptions.
\end{exercise}

\begin{exercise}[1]
Define a contrast.
\end{exercise}

\begin{exercise}[1]
Define an interaction effect.
\end{exercise}

\begin{exercise}[2]
Suppose \(k=4\) and \(N=40\). Find the treatment, error, and total
degrees of freedom.
\end{exercise}

\begin{exercise}[2]
Suppose

\[
SS_A=90,
\qquad
SS_E=210,
\qquad
k=4,
\qquad
N=32.
\]

Compute \(MS_A\), \(MS_E\), and \(F\).
\end{exercise}

\begin{exercise}[2]
Suppose

\[
SS_T=500,
\qquad
SS_E=350.
\]

Find \(SS_A\).
\end{exercise}

\begin{exercise}[2]
Three groups have sizes

\[
5,\ 10,\ 15
\]

and means

\[
12,\ 10,\ 8.
\]

Compute the grand mean.
\end{exercise}

\begin{exercise}[2]
Suppose

\[
SS_A=120,
\qquad
SS_T=300.
\]

Compute \(\eta^2\).
\end{exercise}

\begin{exercise}[2]
Suppose

\[
SS_A=80,
\qquad
SS_E=120.
\]

Compute partial eta-squared.
\end{exercise}

\begin{exercise}[2]
How many pairwise comparisons are possible among \(6\) groups?
\end{exercise}

\begin{exercise}[2]
If \(15\) comparisons are Bonferroni-adjusted to control FWER at \(0.05\),
find the individual significance level.
\end{exercise}

\begin{exercise}[2]
Construct a contrast comparing group \(1\) with the average of groups
\(2,3,4\).
\end{exercise}

\begin{exercise}[2]
Determine whether the coefficient vector

\[
(1,-1,0)
\]

defines a contrast.
\end{exercise}

\begin{exercise}[2]
Determine whether

\[
(1,1,-2)
\]

defines a contrast.
\end{exercise}

\begin{exercise}[2]
A \(2\times3\) factorial experiment has \(4\) replicates per cell.
Find the total sample size.
\end{exercise}

\begin{exercise}[2]
For the previous design, compute the degrees of freedom for factor \(A\),
factor \(B\), interaction, error, and total.
\end{exercise}

\begin{exercise}[3]
Derive the identity

\[
SS_T
=
SS_A+SS_E.
\]
\end{exercise}

\begin{exercise}[3]
Explain why the cross term vanishes in the ANOVA decomposition.
\end{exercise}

\begin{exercise}[3]
Show that

\[
df_T
=
df_A+df_E.
\]
\end{exercise}

\begin{exercise}[3]
Explain why \(MS_E\) estimates the common within-group variance.
\end{exercise}

\begin{exercise}[3]
Explain why \(MS_A\) tends to become large when population means differ.
\end{exercise}

\begin{exercise}[3]
Show that one-way ANOVA can be written as a regression model with
indicator variables.
\end{exercise}

\begin{exercise}[3]
Show that the one-way ANOVA null is equivalent to jointly testing all
group-indicator coefficients equal to zero.
\end{exercise}

\begin{exercise}[3]
Prove that for two groups, the pooled two-sample test satisfies

\[
F=t^2.
\]
\end{exercise}

\begin{exercise}[3]
Derive the standard error of a planned contrast.
\end{exercise}

\begin{exercise}[3]
Construct two orthogonal contrasts among four equally sized groups.
\end{exercise}

\begin{exercise}[3]
Show that at most \(k-1\) linearly independent contrasts exist among
\(k\) group means.
\end{exercise}

\begin{exercise}[3]
Explain why multiple unadjusted pairwise tests inflate family-wise Type I
error.
\end{exercise}

\begin{exercise}[3]
Compare Bonferroni and Tukey methods conceptually.
\end{exercise}

\begin{exercise}[3]
Explain why unequal variances combined with unequal group sizes can be
problematic for ordinary ANOVA.
\end{exercise}

\begin{exercise}[3]
Explain why Kruskal--Wallis is not automatically a test of medians.
\end{exercise}

\begin{exercise}[3]
For a two-way model, explain the difference between a main effect and an
interaction.
\end{exercise}

\begin{exercise}[3]
Construct a \(2\times2\) table of means that has no interaction.
\end{exercise}

\begin{exercise}[3]
Construct a \(2\times2\) table of means that has strong crossover
interaction.
\end{exercise}

\begin{exercise}[3]
Explain why strong interaction can make marginal main effects
misleading.
\end{exercise}

\begin{exercise}[3]
Explain how blocking can reduce residual variability.
\end{exercise}

\begin{exercise}[3]
Distinguish crossed and nested factors.
\end{exercise}

\begin{exercise}[3]
Explain how ANCOVA combines ANOVA and regression.
\end{exercise}

\begin{exercise}[3]
Explain the meaning of a random effect.
\end{exercise}

\begin{exercise}[3]
Derive the intraclass correlation

\[
\rho
=
\frac{
\sigma_A^2
}{
\sigma_A^2+\sigma^2
}
\]

for a random-intercept model.
\end{exercise}

\begin{exercise}[3]
Explain why unbalanced factorial ANOVA can require careful choice of sums
of squares.
\end{exercise}

\begin{exercise}[4]
Derive the exact \(F\)-distribution of the one-way ANOVA statistic under
normal equal-variance assumptions.
\end{exercise}

\begin{exercise}[4]
Derive the expected mean squares under the one-way fixed-effects model.
\end{exercise}

\begin{exercise}[4]
Derive the one-way ANOVA model in matrix form.
\end{exercise}

\begin{exercise}[4]
Show how ANOVA contrasts can be written as general linear hypotheses

\[
R\boldsymbol{\beta}
=
\mathbf{r}.
\]
\end{exercise}

\begin{exercise}[4]
Derive the general partial \(F\)-test for nested linear models.
\end{exercise}

\begin{exercise}[4]
Study Scheff\'e simultaneous confidence intervals.
\end{exercise}

\begin{exercise}[4]
Derive Tukey's HSD procedure using the studentized range distribution.
\end{exercise}

\begin{exercise}[4]
Study the Tukey--Kramer method for unequal sample sizes.
\end{exercise}

\begin{exercise}[4]
Investigate Dunnett's procedure for comparing treatments with a control.
\end{exercise}

\begin{exercise}[4]
Study Welch's ANOVA and derive its approximate degrees of freedom.
\end{exercise}

\begin{exercise}[4]
Compare Bartlett, Levene, and Brown--Forsythe tests.
\end{exercise}

\begin{exercise}[4]
Study the Kruskal--Wallis rank test.
\end{exercise}

\begin{exercise}[4]
Derive the balanced two-way ANOVA sum-of-squares decomposition.
\end{exercise}

\begin{exercise}[4]
Derive expected mean squares for two-way fixed and random-effects models.
\end{exercise}

\begin{exercise}[4]
Study randomized complete block designs.
\end{exercise}

\begin{exercise}[4]
Study repeated-measures ANOVA and sphericity.
\end{exercise}

\begin{exercise}[4]
Investigate Greenhouse--Geisser corrections.
\end{exercise}

\begin{exercise}[4]
Study nested ANOVA models.
\end{exercise}

\begin{exercise}[4]
Study mixed-effects models as generalizations of ANOVA.
\end{exercise}

\begin{exercise}[4]
Investigate Type I, Type II, and Type III sums of squares in unbalanced
designs.
\end{exercise}

\begin{exercise}[4]
Study ANCOVA with treatment-by-covariate interactions.
\end{exercise}

%%%%%%%%%%%%%%%%%%%%%%%%%%%%%%%%%%%%%%%%%%%%%%%%%%%%%%%%%%%%
\subsection{Conceptual Summary}
\label{subsec:anova-summary}
%%%%%%%%%%%%%%%%%%%%%%%%%%%%%%%%%%%%%%%%%%%%%%%%%%%%%%%%%%%%

One-way ANOVA compares

\[
\boxed{
H_0:
\mu_1
=
\mu_2
=
\cdots
=
\mu_k
}
\]

against the alternative that not all means are equal.

The total sample size is

\[
\boxed{
N
=
\sum_{i=1}^{k}
n_i.
}
\]

The grand mean is

\[
\boxed{
\overline{Y}_{\cdot\cdot}
=
\frac{
\sum_i
n_i\overline{Y}_{i\cdot}
}{
N
}.
}
\]

Total variability is

\[
\boxed{
SS_T
=
\sum_{i,j}
\left(
Y_{ij}
-
\overline{Y}_{\cdot\cdot}
\right)^2.
}
\]

Between-group variability is

\[
\boxed{
SS_A
=
\sum_i
n_i
\left(
\overline{Y}_{i\cdot}
-
\overline{Y}_{\cdot\cdot}
\right)^2.
}
\]

Within-group variability is

\[
\boxed{
SS_E
=
\sum_{i,j}
\left(
Y_{ij}
-
\overline{Y}_{i\cdot}
\right)^2.
}
\]

These satisfy

\[
\boxed{
SS_T
=
SS_A
+
SS_E.
}
\]

The corresponding mean squares are

\[
\boxed{
MS_A
=
\frac{
SS_A
}{
k-1
},
}
\]

and

\[
\boxed{
MS_E
=
\frac{
SS_E
}{
N-k
}.
}
\]

The ANOVA statistic is

\[
\boxed{
F
=
\frac{
MS_A
}{
MS_E
}.
}
\]

Under the classical null,

\[
\boxed{
F
\sim
F_{k-1,N-k}.
}
\]

ANOVA extends naturally to:

\[
\boxed{
\text{contrasts},
}
\]

\[
\boxed{
\text{multiple comparisons},
}
\]

\[
\boxed{
\text{factorial designs},
}
\]

\[
\boxed{
\text{blocking},
}
\]

\[
\boxed{
\text{repeated measures},
}
\]

and

\[
\boxed{
\text{mixed models}.
}
\]

%%%%%%%%%%%%%%%%%%%%%%%%%%%%%%%%%%%%%%%%%%%%%%%%%%%%%%%%%%%%
\subsection{Section Connections}
\label{subsec:anova-connections}
%%%%%%%%%%%%%%%%%%%%%%%%%%%%%%%%%%%%%%%%%%%%%%%%%%%%%%%%%%%%

\chapterconnection{
Analysis of Variance is a direct extension of the regression framework
from Section~\ref{sec:regression}. Categorical group indicators create a
linear model, and the ANOVA \(F\)-test is a general linear-model test of
several coefficients simultaneously. The decomposition
\(SS_T=SS_A+SS_E\) is the categorical-predictor version of the regression
decomposition \(SST=SSR+SSE\). The next section, Experimental Design,
develops the principles that determine how treatments, factors, blocks,
replicates, and randomization should be arranged before ANOVA is applied.
Later sections use ANOVA ideas in multivariate statistics, time series,
statistical learning, mixed models, and generalized linear modeling.
}

%%%%%%%%%%%%%%%%%%%%%%%%%%%%%%%%%%%%%%%%%%%%%%%%%%%%%%%%%%%%
% SECTION SOURCE TRACKING
%%%%%%%%%%%%%%%%%%%%%%%%%%%%%%%%%%%%%%%%%%%%%%%%%%%%%%%%%%%%

%------------------------------------------------------------
% 8.8 Analysis of Variance
%------------------------------------------------------------
%
% Source basis:
%   - Standard one-way and factorial ANOVA theory.
%   - Standard linear-model theory.
%   - Standard multiple-comparison and experimental-design concepts.
%
% Used for:
%   - one-way ANOVA;
%   - group and grand means;
%   - total, treatment, and error sums of squares;
%   - ANOVA decomposition;
%   - degrees of freedom;
%   - mean squares;
%   - F-statistics;
%   - ANOVA tables;
%   - model assumptions;
%   - residual diagnostics;
%   - ANOVA-regression equivalence;
%   - two-group F=t^2 equivalence;
%   - planned contrasts;
%   - orthogonal contrasts;
%   - pairwise comparisons;
%   - Bonferroni adjustment;
%   - Tukey and Tukey--Kramer previews;
%   - effect sizes;
%   - eta-squared;
%   - partial eta-squared;
%   - omega-squared preview;
%   - Welch ANOVA;
%   - variance-homogeneity procedures;
%   - Kruskal--Wallis preview;
%   - two-way ANOVA;
%   - main effects;
%   - interactions;
%   - factorial designs;
%   - blocking;
%   - randomized complete block designs;
%   - repeated-measures ANOVA;
%   - sphericity preview;
%   - mixed-model preview;
%   - fixed and random effects;
%   - variance components;
%   - intraclass correlation;
%   - nested factors;
%   - ANCOVA;
%   - unbalanced designs;
%   - Type I, II, and III sums of squares;
%   - general linear hypotheses;
%   - partial F-tests;
%   - power and effect-size previews;
%   - applications and exercises.
%
% Notes:
%   - Experimental Design is developed next in Section 8.9.
%   - Regression in Section 8.7 provides the general linear-model
%     framework underlying ANOVA.
%   - Multivariate Statistics later extends ANOVA to multiple response
%     variables.
%   - Statistical Learning and mixed-model methods generalize these ideas
%     beyond classical balanced designs.
%
%%%%%%%%%%%%%%%%%%%%%%%%%%%%%%%%%%%%%%%%%%%%%%%%%%%%%%%%%%%%